%% LyX 2.4.4 created this file.  For more info, see https://www.lyx.org/.
%% Do not edit unless you really know what you are doing.
\documentclass[english]{article}
\usepackage[T1]{fontenc}
\usepackage[utf8]{inputenc}
\setcounter{secnumdepth}{-2}
\setlength{\parindent}{0bp}
\usepackage{amsmath}
\usepackage{amsthm}
\usepackage{amssymb}

\makeatletter
%%%%%%%%%%%%%%%%%%%%%%%%%%%%%% Textclass specific LaTeX commands.
\theoremstyle{plain}
\newtheorem*{thm*}{\protect\theoremname}

\@ifundefined{date}{}{\date{}}
%%%%%%%%%%%%%%%%%%%%%%%%%%%%%% User specified LaTeX commands.
\usepackage{graphicx}
\usepackage{geometry}
\newcommand{\limi}{\underset{n\rightarrow\infty}{\lim}}
\newcommand{\limxz}{\underset{x\rightarrow x_{0}}{\lim}}
\newcommand{\limfxz}{\underset{x\rightarrow x_{0}}{\lim}f(x)}
\newcommand{\limfxm}{\underset{x\rightarrow x_{0}^{-}}{\lim}f(x)}
\newcommand{\limfxp}{\underset{x\rightarrow x_{0}^{+}}{\lim}f(x)}
\newcommand{\limxm}{\underset{x\rightarrow x_{0}^{-}}{\lim}}
\newcommand{\limxp}{\underset{x\rightarrow x_{0}^{+}}{\lim}}
\newcommand{\sqrm}{M_{m\times m}(\mathbb{F})}
\newcommand{\seqa}{(a_{n})_{n=1}^{\infty}}
\newcommand{\seqx}{(x_{n})_{n=1}^{\infty}}
\newcommand{\funcdr}{f:D\rightarrow\mathbb{R}}
\newcommand{\der}{\limxz\frac{f(x)-f(x_{0})}{x-x_{0}}}
\usepackage[all]{pdfprivacy}

\makeatother

\usepackage{babel}
\providecommand{\theoremname}{Theorem}

\begin{document}
\title{{\Large Weierstrass boundedness theorem\vspace{-2em}}}
\maketitle
\begin{thm*}
Let $f:\left[a,b\right]\rightarrow\mathbb{R}$ be a continuous function
from $\left[a,b\right]$ to $\mathbb{R}$.

Then $f$ is bounded on that interval. 
\end{thm*}
\medskip{}

\rule[0.5ex]{1\columnwidth}{1pt}

\medskip{}

\begin{proof}
We will first prove that $f$ is bounded from above. 

Let us assume for the sake of contradiction that $f$ is not bounded
from above. Then 
\[
\forall M\in\mathbb{R}\ \exists x\in\left[a,b\right]\ s.t.\ f(x)>M
\]
and in particular, since $\mathbb{N}\subseteq\mathbb{R}$,

\[
\forall n\in\mathbb{N}\ \exists x_{n}\in\left[a,b\right]\ s.t.\ f(x_{n})>n
\]

Let us look at the sequence $\left(x_{n}\right)_{n=1}^{\infty}$.
Note that $\forall n\in\mathbb{N},a\le x_{n}\le b$, i.e. this is
a bounded sequence.

Thus, the Bolzano-Weierstrass theorem states it has a convergent subsequence,
$\left(x_{n_{k}}\right)_{k=1}^{\infty}$.

Since limits of sequences weakly preserve inequalities maintained
by all elements of the sequence,
\[
a\le\underset{k\rightarrow\infty}{\lim}x_{n_{k}}\le b
\]
This gives $\underset{k\rightarrow\infty}{\lim}x_{n_{k}}=x_{0}\in\left[a,b\right]\subseteq\mathbb{R}$.

Therefore, by Heine's characterization of continuous functions, 

\[
\underset{k\rightarrow\infty}{\lim}f\left(x_{n_{k}}\right)=f(x_{0})\in\mathbb{R}
\]

However, notice that 
\[
\forall k\in\mathbb{N},\ f(x_{n_{k}})>n_{k}\ge k
\]

hence
\[
\underset{k\rightarrow\infty}{\lim}f\left(x_{n_{k}}\right)=\infty
\]

This is a contradiction. This means that $f$ is bounded from above.

To show that $f$ is bounded from below, we can apply the same argument
to $-f$ since $-f$ is also continuous as a scalar multiple of a
continuous function.

Thus, we get the same contradiction again, hence $f$ is bounded.
\end{proof}

\end{document}
